\documentclass[10pt,a4paper]{article}
\usepackage[utf8]{inputenc}
\usepackage[T1]{fontenc}
\usepackage[spanish,es-nodecimaldot]{babel}
\usepackage{lmodern}
\usepackage[a4paper,margin=1.75cm,headheight=14pt]{geometry}
\usepackage{graphicx}
\usepackage{booktabs}
\usepackage{tabularx}
\usepackage{array}
\usepackage{float}
\usepackage{xcolor}
\usepackage{fancyhdr}
\usepackage{microtype}
\usepackage{hyperref}
\usepackage{enumitem}
\usepackage{titlesec}

\definecolor{unabblue}{HTML}{123B63}
\definecolor{softgray}{HTML}{F2F4F6}
\hypersetup{colorlinks=true,linkcolor=unabblue,urlcolor=unabblue,citecolor=unabblue}
\setlength{\parindent}{0pt}
\setlength{\parskip}{4pt}
\setlist{nosep,leftmargin=1.2em}
\titleformat{\section}{\large\bfseries\color{unabblue}}{\thesection.}{0.5em}{}
\titleformat{\subsection}{\normalsize\bfseries\color{unabblue}}{\thesubsection.}{0.5em}{}
\pagestyle{fancy}
\fancyhf{}
\lhead{MCDI501 - Estadística Computacional}
\rhead{Sumativa 2}
\cfoot{\thepage}
\newcommand{\figdir}{../figuras}
\newcommand{\compactcaption}[1]{\vspace{-3pt}\begin{center}\footnotesize #1\end{center}\vspace{-5pt}}

\begin{document}

% Página 1: portada
\thispagestyle{empty}
\begin{center}
\vspace*{1.2cm}
\includegraphics[width=5cm]{logo_unab.png}

\vspace{1.5cm}
{\Large\bfseries MAGÍSTER EN CIENCIA DE DATOS}\par
\vspace{0.5cm}
{\large MCDI501 - Estadística Computacional}\par
\vspace{1.8cm}
{\Huge\bfseries\color{unabblue} Evaluación Sumativa 2}\par
\vspace{0.5cm}
{\LARGE Validación, simulación y remuestreo}\par
\vspace{0.9cm}
{\Large Análisis del riesgo de enfermedad coronaria a diez años}\par

\vfill
\renewcommand{\arraystretch}{1.35}
\begin{tabular}{>{\bfseries}r l}
Estudiante: & Alexander Sepúlveda Ormazábal \\
Docente: & Jean Paul Maidana González \\
Conjunto de datos: & Framingham Heart Study \\
Fecha: & 23 de octubre de 2026 \\
Repositorio: & \href{https://github.com/alexander-seor/MCDI501-Proyecto-Framingham}{MCDI501-Proyecto-Framingham}
\end{tabular}
\vfill
\end{center}
\newpage

% Página 2: propósito y diseño
\section{Propósito y estrategia de validación}

En la Sumativa 1 se estimaron medias, correlaciones y diferencias entre
personas con y sin enfermedad coronaria. En esta segunda etapa se revisa
cuánto cambian esos resultados al remuestrear los mismos 4.238 registros y
al reducir la influencia de valores extremos. Para ello se aplican
bootstrap, una prueba de permutación, una simulación Monte Carlo y análisis
de sensibilidad. Los datos originales no se modifican.

\begin{table}[H]
\centering\small
\renewcommand{\arraystretch}{1.18}
\begin{tabularx}{\textwidth}{>{\bfseries}p{3.2cm} X p{3.0cm}}
\toprule
Componente & Procedimiento & Configuración \\
\midrule
Medias & Intervalos clásico, bootstrap percentil y BCa & 10.000 remuestras \\
Grupos CHD & Welch de S1 y permutación bilateral & 10.000 permutaciones \\
Correlaciones & Bootstrap pareado de cinco correlaciones punto-biseriales & 3.000 remuestras \\
Monte Carlo & Cohortes generadas desde prevalencia de hipertensión y tasas condicionales de CHD & 10.000 cohortes de 1.000 \\
Robustez & Recorte, winsorización, jackknife y bootstrap de diferencia recortada & Semilla 42 \\
\bottomrule
\end{tabularx}
\caption{Diseño de validación reproducible.}
\end{table}

El intervalo clásico aproxima la distribución muestral de la media; el
percentil toma cuantiles de la distribución bootstrap; y BCa corrige sesgo y
aceleración \cite{efron1987,efron1993}. La permutación usa la
intercambiabilidad de las etiquetas bajo la hipótesis nula \cite{good2005}.
La simulación Monte Carlo aproxima la distribución de resultados hipotéticos
bajo parámetros fijos \cite{robert2004}. El análisis se implementó con
Python, NumPy, pandas, Matplotlib, seaborn y SciPy \cite{virtanen2020}.

\subsection{Criterio de interpretación}

En este informe, un resultado se considera estable cuando cambia poco al
repetir el análisis con remuestras de los mismos datos. Esto corresponde a
una validación interna: no asegura que el resultado se mantenga en otra
población, ni que la relación observada sea causal o tenga por sí sola
utilidad predictiva.

\vfill
\textbf{Fuente de datos.} El archivo analizado corresponde a la versión
publicada por Dileep en Kaggle. La documentación del NHLBI se utilizó como
referencia contextual sobre el Framingham Heart Study \cite{dileep,nhlbi}.
\newpage

% Página 3: bootstrap + permutación
\section{Intervalos bootstrap para las medias}

\begin{table}[H]
\centering\scriptsize
\renewcommand{\arraystretch}{1.12}
\begin{tabular}{lrrrrrrr}
\toprule
Variable & Media & \multicolumn{2}{c}{Clásico} & \multicolumn{2}{c}{Percentil} & \multicolumn{2}{c}{BCa} \\
 & & Inf. & Sup. & Inf. & Sup. & Inf. & Sup. \\
\midrule
Edad & 49,5849 & 49,3268 & 49,8431 & 49,3259 & 49,8462 & 49,3268 & 49,8471 \\
Presión sistólica & 132,3524 & 131,6887 & 133,0161 & 131,6881 & 133,0109 & 131,7004 & 133,0231 \\
Colesterol total & 236,7216 & 235,3707 & 238,0724 & 235,3656 & 238,0499 & 235,4043 & 238,1034 \\
\bottomrule
\end{tabular}
\caption{Comparación de intervalos de confianza al 95\%.}
\end{table}

\noindent\makebox[\textwidth][c]{%
\includegraphics[width=\textwidth,height=5.1cm,keepaspectratio]{\figdir/figura_1_bootstrap_intervalos.png}}
\compactcaption{Figura 1. Distribuciones de las medias bootstrap y comparación de intervalos.}

Los tres intervalos entregan prácticamente los mismos límites. Por eso, la
estimación de estas medias cambia poco con el método elegido. De todas
formas, esta coincidencia no resuelve los problemas de datos faltantes,
valores extremos o representatividad de la muestra.

\section{Prueba de permutación}

La presión sistólica media fue 143,62 mmHg con CHD y 130,33 mmHg sin CHD,
una diferencia de 13,29 mmHg. Ninguna de las 10.000 permutaciones produjo
un valor tan extremo, por lo que se obtuvo $p=0{,}00010$. El resultado apunta
en la misma dirección que Welch en la Sumativa 1
($p=1{,}210\times 10^{-30}$), aunque ambas pruebas no son idénticas cuando
las dispersiones de los grupos difieren. En esta muestra, la presión media
fue mayor entre quienes desarrollaron CHD. La comparación no está ajustada
por edad, sexo u otros posibles factores de confusión.
\newpage

% Página 4: permutación y correlaciones
\noindent\makebox[\textwidth][c]{%
\includegraphics[width=\textwidth,height=5.6cm,keepaspectratio]{\figdir/figura_2_permutacion_presion_sistolica.png}}
\compactcaption{Figura 2. Distribución nula y diferencia observada de presión sistólica.}

\section{Estabilidad de las correlaciones}

\begin{table}[H]
\centering\small
\begin{tabular}{lrrrr}
\toprule
Variable & $n$ efectivo & $r$ observado & IC 95\% inferior & IC 95\% superior \\
\midrule
Edad & 4.238 & 0,2253 & 0,1964 & 0,2540 \\
Presión sistólica & 4.238 & 0,2164 & 0,1826 & 0,2494 \\
Presión diastólica & 4.238 & 0,1453 & 0,1115 & 0,1797 \\
Glucosa & 3.850 & 0,1255 & 0,0815 & 0,1680 \\
Colesterol total & 4.188 & 0,0822 & 0,0505 & 0,1135 \\
\bottomrule
\end{tabular}
\caption{Correlaciones con \texttt{TenYearCHD} e intervalos bootstrap.}
\end{table}

\noindent\makebox[\textwidth][c]{%
\includegraphics[width=\textwidth,height=6.7cm,keepaspectratio]{\figdir/figura_3_estabilidad_correlaciones.png}}
\compactcaption{Figura 3. Estimaciones puntuales e intervalos bootstrap al 95\%.}

En las 3.000 remuestras, los intervalos de las cinco correlaciones quedaron
por encima de cero. Esto indica que el signo positivo se mantuvo dentro de
esta muestra. Sin embargo, las asociaciones siguen siendo débiles o
modestas y no bastan, por sí solas, para afirmar que las variables tengan
buena capacidad predictiva. Glucosa, además, tiene menos casos completos y
fue la variable más sensible a los valores extremos.
\newpage

% Página 5: Monte Carlo
\section{Simulación Monte Carlo}

La simulación utiliza $P(H)=0{,}3105$, $P(CHD\mid H)=0{,}2470$ y
$P(CHD\mid \neg H)=0{,}1092$. Para cada cohorte se genera primero la
hipertensión y luego CHD condicionado a ese estado. Los parámetros permanecen
fijos y las personas se consideran independientes.

\begin{table}[H]
\centering\small
\begin{tabular}{lr}
\toprule
Resultado & Valor \\
\midrule
Tasa global observada y teórica & 15,20\% \\
Media de las 10.000 cohortes & 15,19\% \\
Desviación entre cohortes & 1,13 puntos porcentuales \\
MCSE de la media simulada & 0,0113 puntos porcentuales \\
Rango predictivo central 95\% & 13,0\%--17,4\% \\
Percentil 99 & 17,9\% \\
$P(\mathrm{tasa}\geq 18\%)$ & 0,88\% \\
\bottomrule
\end{tabular}
\caption{Resumen de 10.000 cohortes simuladas de 1.000 personas.}
\end{table}

\noindent\makebox[\textwidth][c]{%
\includegraphics[width=\textwidth,height=8.1cm,keepaspectratio]{\figdir/figura_4_monte_carlo_chd.png}}
\compactcaption{Figura 4. Distribución de tasas simuladas y convergencia de la media.}

La media de las cohortes simuladas fue 15,19\%, muy cercana al 15,20\%
esperado a partir de los parámetros usados. Esto comprueba que la simulación
fue implementada de manera coherente, pero no constituye una validación en
otra población porque las probabilidades provienen de los mismos datos. El
rango 13,0\%--17,4\% describe variación entre cohortes simuladas, no un
intervalo de confianza poblacional. El umbral de 18\% es analítico y no un
punto de corte clínico.
\newpage

% Página 6: robustez y S3
\section{Análisis de robustez}

\begin{table}[H]
\centering\scriptsize
\begin{tabular}{lrrrrrr}
\toprule
Variable & Media & Mediana & Recortada 5\% & Winsorizada & EE jack. & Extremo influyente \\
\midrule
Edad & 49,585 & 49,0 & 49,447 & 49,583 & 0,1317 & 70 \\
Presión sistólica & 132,352 & 128,0 & 130,906 & 132,240 & 0,3385 & 295 \\
Colesterol total & 236,722 & 234,0 & 235,193 & 236,407 & 0,6890 & 696 \\
\bottomrule
\end{tabular}
\caption{Sensibilidad de tendencia central e influencia jackknife.}
\end{table}

La diferencia de presión sistólica cambió poco al limitar los valores
extremos: pasó de 13,29 a 12,78 mmHg. Con medias recortadas fue de 13,15
mmHg y su intervalo bootstrap, de 10,99 a 15,35 mmHg, permaneció por encima
de cero. Por tanto, la diferencia no parece depender de unas pocas
observaciones extremas. Las correlaciones también conservaron su signo,
aunque glucosa mostró el mayor cambio, de 0,1255 a 0,1088, y debe
interpretarse con más cautela.

\section{Preparación para la Sumativa 3}

\begin{table}[H]
\centering\scriptsize
\renewcommand{\arraystretch}{1.12}
\begin{tabularx}{\textwidth}{>{\bfseries}p{2.7cm} p{3.7cm} p{2.7cm} X}
\toprule
Resultado & Evidencia de S2 & Estado & Recomendación \\
\midrule
Medias principales & IC clásico, percentil y BCa coinciden & Robusto & Controlar extremos dentro del entrenamiento. \\
Diferencia de sysBP & Permutación, winsorización y recorte coinciden & Robusto como asociación & Ajustar confusores; no inferir causalidad. \\
Correlaciones & $r=0{,}082$--$0{,}225$; IC excluyen cero & Estable pero débil & No elegir variables solo por correlación marginal. \\
Glucosa & $r$: 0,1255 a 0,1088; $n=3.850$ & Usar con cautela & Estudiar faltantes y extremos. \\
Tasa de CHD & 15,20\% observada; 15,19\% simulada & Coherente bajo supuestos & Estratificar y validar fuera de muestra. \\
Calidad de datos & sysBP=295; colesterol=696; faltantes & Revisar & Verificar plausibilidad e imputar sin fuga. \\
\bottomrule
\end{tabularx}
\caption{Síntesis para la etapa predictiva.}
\end{table}

Antes de imputar, escalar, transformar o winsorizar se deben separar las
particiones; cada operación se ajustará solo con entrenamiento. Dado que CHD
ocurre en 15,2\%, se recomiendan particiones estratificadas y evaluación de
discriminación y calibración fuera de muestra.
\newpage

% Página 7: conclusiones y referencias
\section{Conclusiones}

Ninguna de las comprobaciones realizadas cambió las conclusiones principales
de la Sumativa 1. La edad y la presión sistólica mostraron los resultados más
estables dentro de esta muestra, por lo que tiene sentido considerarlas en la
siguiente etapa. Las correlaciones conservaron el signo, pero fueron pequeñas
y no demuestran capacidad predictiva por sí solas. Glucosa requiere más
cautela por sus datos faltantes y su mayor sensibilidad a los extremos. Estas
conclusiones orientan la modelación, pero todavía deben comprobarse fuera de
muestra y no permiten establecer causalidad.

\small
\begin{thebibliography}{9}
\bibitem{dileep} Dileep. (2019). \textit{Logistic regression to predict heart disease} [Conjunto de datos]. Kaggle. \url{https://www.kaggle.com/datasets/dileep070/heart-disease-prediction-using-logistic-regression}
\bibitem{efron1987} Efron, B. (1987). Better bootstrap confidence intervals. \textit{Journal of the American Statistical Association, 82}(397), 171--185. \url{https://doi.org/10.1080/01621459.1987.10478410}
\bibitem{efron1993} Efron, B., \& Tibshirani, R. J. (1993). \textit{An introduction to the bootstrap}. Chapman \& Hall.
\bibitem{good2005} Good, P. I. (2005). \textit{Permutation, parametric, and bootstrap tests of hypotheses} (3.a ed.). Springer.
\bibitem{nhlbi} National Heart, Lung, and Blood Institute. (s. f.). \textit{Teaching datasets: Framingham Heart Study}. \url{https://biolincc.nhlbi.nih.gov/teaching/}
\bibitem{robert2004} Robert, C. P., \& Casella, G. (2004). \textit{Monte Carlo statistical methods} (2.a ed.). Springer.
\bibitem{virtanen2020} Virtanen, P., et al. (2020). SciPy 1.0: Fundamental algorithms for scientific computing in Python. \textit{Nature Methods, 17}(3), 261--272. \url{https://doi.org/10.1038/s41592-019-0686-2}
\end{thebibliography}

\vfill
\begin{center}\footnotesize
Notebook reproducible ejecutado con semilla 42. Todas las figuras fueron
exportadas a 300 dpi y los datos originales permanecieron sin cambios.
\end{center}

\end{document}
